\documentclass[11pt,a4paper]{article}
\usepackage[utf8]{inputenc}
\usepackage[T1]{fontenc}
\usepackage{lmodern}
\usepackage[margin=1in]{geometry}
\usepackage{enumitem}
\usepackage{xcolor}
\usepackage{hyperref}
\usepackage{booktabs}
\usepackage{longtable}
\usepackage{array}
\usepackage[most]{tcolorbox}

% Define custom colors
\definecolor{promptblue}{RGB}{41, 98, 150}
\definecolor{promptbg}{RGB}{240, 247, 255}
\definecolor{promptborder}{RGB}{100, 160, 220}
\definecolor{accentgreen}{RGB}{0, 128, 100}
\definecolor{warningorange}{RGB}{200, 120, 0}

\newtcolorbox{resultbox}[1][]{
    colback=promptbg,
    colframe=promptborder,
    coltitle=white,
    fonttitle=\bfseries\small,
    title={#1},
    colbacktitle=promptblue,
    boxrule=1pt,
    arc=2mm,
    left=8pt, right=8pt, top=6pt, bottom=6pt,
    enhanced, breakable
}

\newtcolorbox{gapbox}[1][]{
    colback=orange!5,
    colframe=warningorange,
    coltitle=white,
    fonttitle=\bfseries\small,
    title={#1},
    colbacktitle=warningorange,
    boxrule=1pt,
    arc=2mm,
    left=8pt, right=8pt, top=6pt, bottom=6pt,
    enhanced, breakable
}

\hypersetup{
    colorlinks=true,
    linkcolor=promptblue,
    urlcolor=promptblue,
    citecolor=promptblue
}

\title{\textbf{Problem Statement: hardware in the loop GPT }}
\author{Generated by Research Accelerant Agent}
\date{\today}

\begin{document}
\maketitle



\begin{resultbox}[Definition]
In the context of a research proposal, a problem statement is a concise, clear, and evidence-based description of a specific \textbf{researchable} issue, a \textbf{demonstrable} gap in knowledge, or a \textbf{measurable} real-world challenge that your research is \textbf{designed to solve}.
\end{resultbox}

\vspace{1em}

\textbf{\large Structured Problem Statement}\par\vspace{0.5em}

\textbf{1. What is Known}\\
The literature on hardware in the loop GPT  has established foundational knowledge across Across the 5 reviewed studies on hardware in the loop GPT , the following methodological patterns emerged:. Key studies have contributed empirical evidence and theoretical frameworks.

\vspace{0.5em}

\textbf{2. What is Missing}\\
Geographic scope: Most studies on hardware in the loop GPT  may be concentrated in specific regions, limiting generalizability to diverse populations.

\vspace{0.5em}

\textbf{3. Affected Stakeholders}\\
Researchers, practitioners, and policymakers working with hardware in the loop GPT  are directly affected by this gap, as are the populations ultimately served by evidence-based interventions.

\vspace{0.5em}

\textbf{4. Consequences of Inaction}\\
Without addressing this gap, decision-makers will continue to rely on incomplete or non-generalizable evidence, potentially leading to ineffective strategies and missed opportunities for improvement in hardware in the loop GPT .

\vspace{0.5em}

\textbf{5. How the Study Addresses the Gap}\\
The proposed study will employ rigorous methodology to directly investigate geographic scope: most studies on hardware in the loop gpt  may be concentrated in specific regions, limiting generalizability to diverse populations., generating actionable evidence to fill this critical gap in the literature.

\vspace{1em}

\begin{resultbox}[Full Problem Statement]
Despite advances in understanding hardware in the loop GPT , a critical gap remains: geographic scope: most studies on hardware in the loop gpt  may be concentrated in specific regions, limiting generalizability to diverse populations.. This is problematic because current evidence is insufficient to guide effective practice across diverse contexts. The proposed study addresses this gap by conducting rigorous, contextually grounded research that will produce generalizable findings and actionable recommendations for stakeholders in the field.
\end{resultbox}

\vspace{1em}

\textbf{\large Validation Checklist}\par\vspace{0.3em}

\begin{itemize}[leftmargin=*, label=$\square$]
    \item Is the problem \textbf{researchable}? (Can it be investigated through empirical methods?)
    \item Is the gap \textbf{demonstrable}? (Can you cite evidence that it exists?)
    \item Is the challenge \textbf{measurable}? (Can outcomes be quantified or assessed?)
    \item Are \textbf{stakeholders} clearly identified?
    \item Are \textbf{consequences of inaction} articulated?
    \item Does the statement tell reviewers \textbf{why}, \textbf{what}, and \textbf{how}?
\end{itemize}

\vfill
\begin{center}
\textit{Problem statement generated by Research Accelerant Agent on \today.}
\end{center}

\end{document}